\section{The certified bounds on \texorpdfstring{$\gamma_5$ and $\gamma_6$}{gamma\_5 and gamma\_6}}\label{app:cert}

This appendix states what the program \texttt{cert.c} of the accompanying repository checks and why its
checks are sound. For a finite set $P$ of $n$ distinct points in the plane, the $L=\binom n2$ pairs give $L$
directions in $\R/\pi\Z$, counted with multiplicity; $d_p$ is the direction of the line through the pair $p$. Fix a cyclic order of the $L$ pairs compatible with their
directions (pairs with equal directions in any order). Consecutive pairs in this order bound $L$ gaps
$g_1,\dots,g_L\ge0$ with $\sum g_i=\pi$, and $G(P)=\max g_i$; $\gamma_n=\inf_P G(P)$. Lower bounds
$\gamma_n\ge\theta$ are proved by contradiction from the assumption
\begin{equation}\label{eq:cert-ass}
  g_i<\theta\quad\text{for all }i,\qquad\text{with}\qquad 0<E:=L\theta-\pi<\theta .
\end{equation}

\begin{lemma}\label{lem:cert-H}
Assume \eqref{eq:cert-ass}. For pairs $p\ne q$, let $\ell(p\to q)$ be the sum of the gaps from $p$ to $q$ in the
cyclic order and $m\in\{1,\dots,L-1\}$ their number, i.e.\ $m-1$ pairs lie strictly between $p$ and $q$. Then
$m\theta-E<\ell(p\to q)<m\theta$. In particular, $\ell(p\to q)>\theta-E>0$.
\end{lemma}

\begin{proof}
$\ell(p\to q)$ is a sum of $m$ gaps, each $<\theta$, and $\pi-\ell(p\to q)$ is a sum of the other $L-m$ gaps,
so $\pi-\ell(p\to q)<(L-m)\theta$.
\end{proof}

Taking $m=1$, every gap exceeds $\theta-E>0$: under \eqref{eq:cert-ass} the $L$ directions are pairwise distinct,
the cyclic order is unique, $\ell(p\to q)$ is the counterclockwise angle from $d_p$ to $d_q$, and $m-1$ is the number
of pairs $r$ whose direction lies in the open counterclockwise arc from $d_p$ to $d_q$.

\paragraph{Normalisation.} Similarities preserve the multiset of gaps (reflections reverse the cyclic
order). So we may assume that a longest pair is $p_0=(0,0)$, $p_1=(1,0)$; then all distances are $\le1$ and
every point lies in $[0,1]\times[-1,1]$. Label the others so that $x_2\le\dots\le x_{n-1}$. The reflection
$(x,y)\mapsto(1-x,y)$, which swaps $p_0,p_1$ and reverses the labels, replaces $x_2+x_{n-1}$ by
$2-(x_2+x_{n-1})$, so we may assume $x_2+x_{n-1}\le1$; then the reflection $(x,y)\mapsto(x,-y)$, which keeps the
$x$-coordinates, gives $y_2\ge0$.

\paragraph{Boxes and direction intervals.} The program covers the normalised configurations,
written as points $(p_2,\dots,p_{n-1})$ of $\R^{2(n-2)}$, by closed axis-parallel boxes with dyadic corners (unit $2^{-40}$) and splits boxes in
half until each is discarded by one of the tests below. For a box and a pair $(i,j)$, every difference
$p_j-p_i$ lies in the Minkowski difference $B_j-B_i$ of the two point boxes. If $B_j-B_i$ does not contain
$0$, the directions of its points form an arc of length $<\pi$ whose endpoints are directions of corners of
$B_j-B_i$; the program encloses each corner direction between two consecutive entries of a table of integer
vectors $v_m\approx2^{62}(\cos\frac{m\pi}N,\sin\frac{m\pi}N)$, $N=921600$, using exact integer cross products, and
widens the result by one table step on each side. This is sound because every table entry has angle within
$1.6\cdot10^{-19}$ of $m\pi/N$, far below one step ($\pi/N\approx3.4\cdot10^{-6}$): the table is generated in high
precision and every entry is checked in interval arithmetic by a separate script, whose output and the table's
SHA-256 hash are recorded with the runs. The program itself checks the two endpoints and that consecutive entries
turn counterclockwise. So each pair gets either
an interval $I_p=[a_p,a_p+w_p]\subset\R/\pi\Z$ containing $d_p$ for every configuration in the box, or the label
\emph{unknown} (when $0\in B_j-B_i$, which includes boxes where two points may coincide, or when the arc is too
long). All thresholds are integers in units of $\pi/N$, and $\theta$ is an integer number of units.

\paragraph{The tests.} A box is discarded if one of the following holds; each implies that no configuration
in the box satisfies the normalisation and \eqref{eq:cert-ass}.
\begin{itemize}
  \item[(D)] The box violates the normalisation: $y_2<0$ on the whole box, or $x_k>x_{k+1}$ on it, or
    $x_2+x_{n-1}>1$ on it, or the boxes of two points are at distance $>1$.
  \item[(P)] Two pairs $p,q$ have known intervals contained in a common arc of length $h\le\theta-E$. Then
    $\ell(p\to q)\le h$ or $\ell(q\to p)\le h$, against Lemma~\ref{lem:cert-H}.
  \item[(J)] Two pairs $p,q$ have known, disjoint intervals with $I_q$ counterclockwise after $I_p$ (within
    the half-turn), so that $\ell(p\to q)$ is the counterclockwise angle and lies in $[\ell^-,\ell^+]$, the
    corresponding interval of angles. The number $m-1$ of pairs strictly between $p$ and $q$ is at least the
    number $j^-$ of pairs whose known interval lies in the open arc from the end of $I_p$ to the start of $I_q$,
    and at most the number $j^+$ of pairs other than $p,q$ that are unknown or whose interval meets the open arc from the start of
    $I_p$ to the end of $I_q$. The test discards the box if no
    $m\in[j^-+1,j^++1]$ has $(m\theta-E,m\theta)\cap[\ell^-,\ell^+]\ne\emptyset$, against
    Lemma~\ref{lem:cert-H}.
  \item[(M)] Fix a pair $p$ with known interval. The map $q\mapsto m_q$ (the number $m$ of
    Lemma~\ref{lem:cert-H} for $\ell(p\to q)$) is a bijection from the other $L-1$ pairs onto $\{1,\dots,L-1\}$.
    For $q$ as in (J), $m_q$ lies in the interval $R_q$ of those $m$ with
    $(m\theta-E,m\theta)\cap[\ell^-,\ell^+]\ne\emptyset$; for every other $q$ put $R_q=\{1,\dots,L-1\}$. The test
    discards the box if there is no bijection $q\mapsto m_q\in R_q$. For intervals $R_q$ this is decided exactly
    (Glover \cite{Glover}) by the greedy rule that assigns $m=1,2,\dots$ in turn to an unused $q$ with $m\in R_q$ and smallest
    $\max R_q$: if some bijection exists, an exchange argument turns it into the greedy one.
\end{itemize}
A box that is not discarded is split at the midpoint of one coordinate. A box whose chosen coordinate has
width at most $2^{-40}$, or whose depth exceeds $200$, is reported as unresolved. The certificate is valid
when the whole initial box $[0,1]^{n-2}\times[-1,1]^{n-2}$ (for $(x_k)$ and $(y_k)$) is processed with no
unresolved box: then no normalised configuration satisfies \eqref{eq:cert-ass}, so $\gamma_n\ge\theta$. In the runs,
the boxes remaining after a breadth-first phase are distributed among several processes with the same
parameters; a run is valid only if every process ends normally with exit status $0$ and reports no unresolved box.

\paragraph{Runs.} With $\theta=64918$ units ($12.6792\ldots^\circ$) for $n=6$ and $\theta=100488$ units
($19.6265\ldots^\circ$) for $n=5$, every box is discarded (45{,}959{,}045 and 1{,}087{,}209 boxes; Table~\ref{tab:gamma}).
Remark~\ref{rem:hexagon} needs only $\gamma_6>12.5495^\circ$, already given by the run with $\theta=64256$ units
($12.55^\circ$, 27{,}631{,}561 boxes). The upper bounds in Table~\ref{tab:gamma} are the gaps of explicit
configurations evaluated in interval arithmetic. Commands, runtimes, outputs and file hashes are in the
repository (\texttt{certificates/\allowbreak gamma6\_\allowbreak gamma5}).
